\documentclass[palette=pink,mode=light]{impaTeXPoster}

\usepackage[utf8]{inputenc}
\usepackage[T1]{fontenc}
\usepackage{lipsum}


\title{Poster Teste}
\subtitle{Exemplo de subtítulo}
\author{Bruno Pereira, Maria Silva, João Souza}
\authoremail{bruno@impa.br, maria@impa.br, joao@impa.br}
\posterlogo{images/impatech.png}    % logo à esquerda
\rightposterlogo{images/tex.png}    % logo à direita

\begin{document}
\maketitle

\useblockstyle{testStyle}


\begin{columns}

\column{0.5}

\block{Introduction}{
    Compactness is one of the central notions in modern mathematics.
    It appears naturally in topology, analysis, geometry, and differential
    equations.
}

\plainblock{ 
\lipsum[1-3]
}

\column{0.5}

\block{Main Theorem}{
Let $K$ be a compact subset of $\mathbb{R}^n$ and
$f:K\to\mathbb{R}$ be continuous. Then $f$ attains a maximum and a
minimum value on $K$.
}

\plainblock{  \textcolor{pink}{\textbf{Proof.}}
Compactness allows one to extract convergent subsequences from bounded
sequences, providing a bridge between local and global properties.

\lipsum[3]
}

\plainblock{ \textcolor{pink}{\textbf{Historical Remark!}}
The modern notion of compactness emerged gradually from the work of
Bolzano, Weierstrass, Borel, and Lebesgue during the nineteenth century.
}

\block{Motivation}{
Many existence theorems require compactness assumptions. For instance,
continuous functions on compact sets attain their extrema.
}

\block{Applications}{
\begin{itemize}
    \item Extreme Value Theorem
    \item Arzelà--Ascoli Theorem
    \item Spectral Theory
    \item Variational Methods
\end{itemize}
}

\end{columns}

\begin{columns}

%Left Column
\column{0.33}

\block{Important Equation}{
The heat equation is given by

\[
\frac{\partial u}{\partial t}
=
\alpha \nabla^2 u.
\]

where $\alpha$ is the thermal diffusivity.\\

\lipsum[3][1-6]
}


%Middle Column
\column{0.33}

\block{Experimental Results}{
\centering

\begin{tabular}{lcc}
\hline
Method & Error & Time (s)\\
\hline
A & 0.014 & 1.2\\
B & 0.008 & 1.8\\
C & 0.005 & 2.4\\
\hline
\end{tabular}
}

\plainblock{\lipsum[1][1-4]}


%Right Column
\column{0.33}

\begin{postercard}{Metodologia}
    \lipsum[7][1-16]
\end{postercard}

\end{columns}

\begin{columns}
\column{0.5}

\plainblock{
\textcolor{pink}{\textbf{Theorem.}}

Let $K \subset \mathbb{R}^n$ be compact and
$f:K\to\mathbb{R}$ continuous.

Then $f$ attains its maximum and minimum values on $K$.

}

\plainblock{
    \posterfigure{0.9}{images/lora.png}%
        {Arquitetura utilizada para o treinamento com LoRA.}{fig:lora}
}

\column{0.5}



\block{Workflow}{
\begin{enumerate}
\item Data acquisition
\item Pre-processing
\item Model training
\item Evaluation
\end{enumerate}
}

\plainblock{
\lipsum[2][1-15]
}

\block{Conclusion}{
We introduced a new framework for studying compactness
and demonstrated several applications in analysis and topology.
}

\end{columns}

\posterfooter{https://github.com/impaTeX/modelo-cartaz/}{%

    \textbf{{\LARGE Referências:}}

    \begin{thebibliography}{99}
    \bibitem{popper}
    POPPER, Karl.
    \textit{A lógica da pesquisa científica}.
    São Paulo: Cultrix, 2002.

    \bibitem{chalmers}
    CHALMERS, A. F.
    \textit{O que é ciência, afinal?}
    São Paulo: Brasiliense, 1999.
    \end{thebibliography}
}

\end{document}